% !TEX root = ../main.tex
\corrigesChapitre{17}{Écoulement des fluides}
\corrigesSouples

\begin{blocLivre}{cor-17-1}
\begin{corrige}{1}{D'où vient la poussée d'Archimède ?}
\demarche{Statique des fluides : à la profondeur $d$, $p=p_\mathrm{atm}+\rho\,g\,d$ ; force pressante $F=p\,S$.}
\begin{questions}
\item $p_\mathrm{haut}=\num{1.013e5}+1000\times9{,}81\times0{,}200=\res{\qty{1,033e5}{Pa}}$ et, à \qty{0,300}{m} de profondeur,
$p_\mathrm{bas}=\res{\qty{1,042e5}{Pa}}$.
\item Avec $S=a^2=\qty{1,00e-2}{m^2}$ : $F_\mathrm{haut}=\res{\qty{1033}{N}}$, vers le bas, et $F_\mathrm{bas}=\res{\qty{1042}{N}}$, vers le
haut.
\item Sur deux faces latérales opposées, la pression est la même à chaque profondeur : les forces sont opposées et de
même valeur. La résultante vaut $F_\mathrm{bas}-F_\mathrm{haut}=(p_\mathrm{bas}-p_\mathrm{haut})\,a^2=\rho\,g\,a^3=\res{\qty{9,81}{N}}$,
vers le haut : c'est exactement $\rho_\mathrm{eau}\,V\,g$.
\item Les deux pressions augmentent de $\rho\,g\times\qty{1}{m}$ : leur différence, donc la résultante, \res{\text{ne change pas}}.
\end{questions}
\end{corrige}
\end{blocLivre}
\begin{blocLivre}{cor-17-2}
\begin{corrige}{2}{La couronne du roi Hiéron}
\demarche{Couronne immobile : $\vec{P}+\vec{\Pi}+\vec{T}=\vec{0}$, la tension $T$ du fil étant l'indication du dynamomètre.}
\begin{questions}
\item Poids $\vec{P}$ vers le bas, poussée $\vec{\Pi}$ et tension $\vec{T}$ vers le haut : $\Pi=P-T=9{,}810-9{,}180=\res{\qty{0,630}{N}}$.
\item $V=\dfrac{\Pi}{\rho_\mathrm{eau}\,g}=\res{\qty{6,42e-5}{m^3}}$, soit \qty{64,2}{cm^3}.
\item $\rho=\dfrac{m}{V}=\res{\qty{1,56e4}{kg.m^{-3}}}$, inférieure à celle de l'or : la couronne \res{\text{n'est pas en or pur}} ;
elle contient un métal moins dense, comme l'argent.
\item $V_\mathrm{or}=\dfrac{1{,}000}{19\,300}=\qty{5,18e-5}{m^3}$, d'où $\Pi=\qty{0,508}{N}$ et une indication de
$\res{\qty{9,30}{N}}$.
\end{questions}
\end{corrige}
\end{blocLivre}
\begin{blocLivre}{cor-17-3}
\begin{corrige}{3}{Le tuyau d'arrosage}
\demarche{$D_V=V/\Delta t=S\,v$, conservé tout au long du tuyau.}
\begin{questions}
\item $D_V=\dfrac{\num{10.0e-3}}{50}=\res{\qty{2,0e-4}{m^3.s^{-1}}}$, soit \res{\qty{12}{L}} par minute.
\item $S=\pi\times(\num{7.5e-3})^2=\qty{1,77e-4}{m^2}$, d'où $v=\dfrac{D_V}{S}=\res{\qty{1,13}{m.s^{-1}}}$.
\item $v'=v\times\left(\dfrac{15}{5{,}0}\right)^2=\res{\qty{10,2}{m.s^{-1}}}$ : la vitesse est multipliée par \res{9}.
\item $h=\dfrac{v'^2}{2\,g}=\res{\qty{5,3}{m}}$. Pincer l'extrémité \res{\text{réduit la section}} : à débit presque inchangé, la
vitesse de sortie augmente et le jet porte plus loin.
\end{questions}
\end{corrige}
\end{blocLivre}
\begin{blocLivre}{cor-17-4}
\begin{corrige}{4}{L'effet Venturi}
\demarche{Conservation du débit, puis relation de Bernoulli à $z$ constant.}
\begin{questions}
\item Dans \res{\text{l'étranglement}} : $v_\mathrm{B}=4\,v_\mathrm{A}=\res{\qty{2,0}{m.s^{-1}}}$.
\item La pression est plus faible dans l'étranglement :
$p_\mathrm{A}-p_\mathrm{B}=\frac12\,\rho\,(v_\mathrm{B}^2-v_\mathrm{A}^2)=500\times(4{,}0-0{,}25)=\res{\qty{1,9e3}{Pa}}$.
\item $\Delta h=\dfrac{p_\mathrm{A}-p_\mathrm{B}}{\rho\,g}=\dfrac{1875}{9810}=\res{\qty{19}{cm}}$.
\item L'air soufflé entre les feuilles est rapide : sa pression est \res{\text{inférieure à la pression atmosphérique}}, qui
s'exerce sur les faces extérieures. Les forces pressantes poussent les feuilles l'une vers l'autre.
\end{questions}
\schema[corrige]{17-venturi}
\end{corrige}
\end{blocLivre}
\begin{blocLivre}{cor-17-5}
\begin{corrige}{5}{Le tonneau percé}
\demarche{Formule de Torricelli, puis mouvement parabolique du chapitre~15 avec une vitesse initiale horizontale.}
\begin{questions}
\item $v=\sqrt{2\times9{,}81\times0{,}80}=\res{\qty{3,96}{m.s^{-1}}}$.
\item $s=\pi\times(\num{5.0e-3})^2=\qty{7,85e-5}{m^2}$ et $D_V=s\,v=\qty{3,11e-4}{m^3.s^{-1}}$, soit \res{\qty{0,31}{L.s^{-1}}} ou
\res{\qty{18,7}{L}} par minute.
\item Chute de \qty{0,50}{m} : $t=\sqrt{\dfrac{2\times0{,}50}{9{,}81}}=\qty{0,319}{s}$, d'où $x=v\,t=\res{\qty{1,26}{m}}$.
\item La hauteur d'eau $h$ diminue, donc la vitesse \res{\text{diminue}} et le jet \res{\text{se rapproche du tonneau}}.
\end{questions}
\end{corrige}
\end{blocLivre}
\begin{blocLivre}{cor-17-6}
\begin{corrige}{6}{Le ballon-sonde}
\demarche{Référentiel terrestre ; vitesse nulle au lâcher, donc pas de frottement : $m\,\vec{a}=\vec{P}+\vec{\Pi}$.}
\begin{questions}
\item Le poids $\vec{P}$, vertical vers le bas, et la poussée d'Archimède de l'air $\vec{\Pi}$, verticale vers le haut.
\item $\Pi=1{,}20\times3{,}0\times9{,}81=\res{\qty{35,3}{N}}$ ; $m=1{,}50+0{,}17\times3{,}0=\qty{2,01}{kg}$, d'où
$P=\res{\qty{19,7}{N}}$.
\item $a=\dfrac{\Pi-P}{m}=\dfrac{15{,}6}{2{,}01}=\res{\qty{7,8}{m.s^{-2}}}$, vers le haut.
\item La pression extérieure diminue : l'hélium se détend et le ballon \res{\text{gonfle}}. L'enveloppe ne pouvant s'étirer
indéfiniment, elle finit par \res{\text{éclater}} ; les instruments redescendent en parachute.
\end{questions}
\end{corrige}
\end{blocLivre}
\begin{blocLivre}{cor-17-7}
\begin{corrige}{7}{Une artère rétrécie}
\demarche{Diamètre divisé par deux : section divisée par quatre, vitesse multipliée par quatre.}
\begin{questions}
\item $S=\pi\times(\num{3.0e-3})^2=\qty{2,83e-5}{m^2}$ et $D_V=S\,v=\res{\qty{1,13e-5}{m^3.s^{-1}}}$, soit
\res{\qty{0,68}{L}} par minute.
\item $v'=4\times0{,}40=\res{\qty{1,6}{m.s^{-1}}}$.
\item $\Delta p=\frac12\times1060\times(1{,}6^2-0{,}40^2)=\res{\qty{1,27e3}{Pa}}$, soit \res{\qty{9,5}{mmHg}}.
\item Si la pression du sang dans la sténose devient \res{\text{inférieure à celle des tissus}} voisins, ceux-ci écrasent la
paroi, ce qui réduit encore la section. Le souffle provient de l'écoulement \res{\text{tourbillonnaire}} du sang très
rapide à la sortie de la sténose.
\end{questions}
\end{corrige}
\end{blocLivre}
\begin{blocLivre}{cor-17-8}
\begin{corrige}{8}{Le débitmètre de Venturi}
\demarche{Statique dans les tubes verticaux, relation de Bernoulli le long de l'axe horizontal de la conduite.}
\begin{questions}
\item $S_\mathrm{A}\,v_\mathrm{A}=S_\mathrm{B}\,v_\mathrm{B}$, d'où $\res{v_\mathrm{B}=\dfrac{S_\mathrm{A}}{S_\mathrm{B}}\,v_\mathrm{A}=4\,v_\mathrm{A}}$.
\item Dans chaque tube, l'eau est au repos : $p_\mathrm{A}=p_\mathrm{atm}+\rho\,g\,h_\mathrm{A}$ et $p_\mathrm{B}=p_\mathrm{atm}+\rho\,g\,h_\mathrm{B}$,
d'où $p_\mathrm{A}-p_\mathrm{B}=\rho\,g\,\Delta h$. Bernoulli donne aussi
$p_\mathrm{A}-p_\mathrm{B}=\frac12\,\rho\,v_\mathrm{A}^2\big((S_\mathrm{A}/S_\mathrm{B})^2-1\big)$ : en égalant, on obtient l'expression de $v_\mathrm{A}$.
\item $v_\mathrm{A}=\sqrt{\dfrac{2\times9{,}81\times0{,}120}{15}}=\res{\qty{0,40}{m.s^{-1}}}$ et
$D_V=S_\mathrm{A}\,v_\mathrm{A}=\qty{7,9e-4}{m^3.s^{-1}}$, soit \res{\qty{0,79}{L.s^{-1}}} ou \res{\qty{2,85}{m^3.h^{-1}}}.
\item $D_V=S_\mathrm{A}\,v_\mathrm{A}$ et $v_\mathrm{A}$ est proportionnelle à $\sqrt{\Delta h}$. Pour un débit double, $\Delta h$ est
multipliée par 4 : $\res{\qty{48}{cm}}$.
\end{questions}
\schema[corrige]{17-venturi}
\end{corrige}
\end{blocLivre}
\begin{blocLivre}{cor-17-9}
\begin{corrige}{9}{Simuler la vidange d'un réservoir}
\demarche{Pendant $\mathrm{d}t$, le volume qui sort par l'orifice est égal au volume perdu par le récipient.}
\begin{questions}
\item La \res{\text{formule de Torricelli}}. Pendant \code{dt}, le volume $s\,v\,\mathrm{d}t$ sort par l'orifice ; la surface
libre baisse de $\dfrac{s}{S}\,v\,\mathrm{d}t$, puisque $S\,|\mathrm{d}h|=s\,v\,\mathrm{d}t$.
\item La boucle s'arrête quand le récipient est vide ($h\leqslant0$) : \code{t} est la \res{\text{durée de vidange}}. La
simulation et l'expression exacte donnent toutes deux \res{\qty{126,2}{s}}.
\item Avant la boucle : \code{temps, hauteurs = [0], [0.20]} ; dans la boucle : \code{temps.append(t)} et
\code{hauteurs.append(h)} ; après : \code{plt.plot(temps, hauteurs)} puis \code{plt.show()}. La courbe décroît de
plus en plus lentement : la surface libre \res{\text{ralentit}} quand $h$ diminue.
\item La section de l'orifice est multipliée par 4 : la durée est \res{\text{divisée par 4}}, environ \res{\qty{32}{s}}.
\item Le jet se \res{\text{contracte}} à la sortie de l'orifice, et les frottements freinent l'eau : le débit réel est
plus faible que prévu, d'environ 17~\% ici.
\end{questions}
\schema[corrige]{17-ex09}
\end{corrige}
\end{blocLivre}
\begin{blocLivre}{cor-17-10}
\begin{corrige}{10}{Le densimètre}
\demarche{À l'équilibre, la poussée compense le poids : $\rho\,S\,h\,g=m\,g$.}
\begin{questions}
\item $\vec{P}+\vec{\Pi}=\vec{0}$, avec $\Pi=\rho\,S\,h\,g$ : $\res{h=\dfrac{m}{\rho\,S}}$. Dans l'eau,
$h=\dfrac{0{,}0240}{1000\times\num{2.00e-4}}=\qty{0,120}{m}$ : on retrouve \qty{12,0}{cm}.
\item $\rho=\dfrac{m}{S\,h}=\dfrac{0{,}0240}{\num{2.00e-4}\times0{,}100}=\res{\qty{1,20e3}{kg.m^{-3}}}$.
\item Il faut $h\leqslant L$, soit $\rho\geqslant\dfrac{m}{S\,L}=\res{\qty{800}{kg.m^{-3}}}$ ; en dessous, le densimètre coule.
\item $h=\res{\qty{10,9}{cm}}$. Les graduations 1000, 1100 et 1200 sont à 12,0 ; 10,9 et 10,0~cm : écarts de 1,1 puis
0,9~cm. Elles \res{\text{ne sont pas régulières}}, car $h$ est inversement proportionnelle à $\rho$.
\end{questions}
\end{corrige}
\end{blocLivre}
\begin{blocLivre}{cor-17-11}
\begin{corrige}{11}{Le château d'eau}
\demarche{Eau immobile : $p=p_\mathrm{atm}+\rho\,g\,h$ ; robinet ouvert : Bernoulli entre la surface libre et la sortie.}
\begin{questions}
\item $p=\num{1.013e5}+1000\times9{,}81\times35=\res{\qty{4,45e5}{Pa}}$, soit environ 4,4~bar.
\item $p=\num{1.013e5}+1000\times9{,}81\times26=\res{\qty{3,56e5}{Pa}}$. La pression est d'autant plus grande que la surface
libre est \res{\text{haute}} : il faut alimenter les derniers étages.
\item Surface libre : $p_\mathrm{atm}$, $v\approx0$ ; sortie : $p_\mathrm{atm}$. D'où $v=\sqrt{2\,g\,h}=\res{\qty{26,2}{m.s^{-1}}}$.
\item $v=\dfrac{D_V}{S}=\dfrac{\num{2.0e-4}}{\num{1.0e-4}}=\res{\qty{2,0}{m.s^{-1}}}$. Les \res{\text{frottements}} dans les
canalisations et le robinet dissipent l'essentiel de l'énergie : la relation de Bernoulli ne s'applique pas ici.
\end{questions}
\end{corrige}
\end{blocLivre}
\begin{blocLivre}{cor-17-12}
\begin{corrige}{12}{Le tube de Pitot}
\demarche{Bernoulli entre un point de l'écoulement libre et le point d'arrêt, à la même altitude.}
\begin{questions}
\item En A, la pression vaut $p_\mathrm{A}$ et la vitesse $v$ ; en B, $v_\mathrm{B}=0$ : $p_\mathrm{A}+\frac12\,\rho\,v^2=p_\mathrm{B}$, d'où
$\res{\Delta p=\frac12\,\rho_\mathrm{air}\,v^2}$.
\item $v=\sqrt{\dfrac{2\,\Delta p}{\rho_\mathrm{air}}}=\sqrt{\dfrac{2\times\num{12.8e3}}{0{,}41}}=\res{\qty{250}{m.s^{-1}}}$, soit
\res{\qty{900}{km.h^{-1}}}.
\item $v=\dfrac{280}{3{,}6}=\qty{77,8}{m.s^{-1}}$ et $\Delta p=\frac12\times1{,}22\times77{,}8^2=\res{\qty{3,7}{kPa}}$.
\item $v_\mathrm{ind}=\sqrt{\dfrac{2\times\num{12.8e3}}{1{,}22}}=\qty{145}{m.s^{-1}}$, environ \res{\qty{520}{km.h^{-1}}}. Cette vitesse
correspond à la même valeur de $\frac12\,\rho\,v^2$ qu'au sol : la portance est la même pour une même vitesse
indiquée, quelle que soit l'altitude. Le pilote garde donc les \res{\text{mêmes repères}}, par exemple la vitesse
minimale de vol.
\end{questions}
\end{corrige}
\end{blocLivre}
\begin{blocLivre}{cor-17-13}
\begin{corrige}{13}{La lance à incendie}
\demarche{Conservation du débit, Bernoulli dans la lance horizontale, puis entre la sortie et le sommet du jet.}
\begin{questions}
\item $D_V=\dfrac{0{,}500}{60}=\qty{8,33e-3}{m^3.s^{-1}}$ ; $v_\mathrm{tuyau}=\dfrac{D_V}{\pi\times0{,}035^2}=\res{\qty{2,17}{m.s^{-1}}}$ et
$v_\mathrm{sortie}=\dfrac{D_V}{\pi\times0{,}011^2}=\res{\qty{21,9}{m.s^{-1}}}$.
\item $p-p_\mathrm{atm}=\frac12\,\rho\,(v_\mathrm{sortie}^2-v_\mathrm{tuyau}^2)=\res{\qty{2,4e5}{Pa}}$, soit environ 2,4~bar.
\item Entre la sortie et le sommet, à l'air libre, où $v=0$ : $h=\dfrac{v_\mathrm{sortie}^2}{2\,g}=\res{\qty{24,5}{m}}$. Atteindre
\qty{20}{m} demande seulement $v=\sqrt{2\,g\times20}=\qty{19,8}{m.s^{-1}}$ : \res{\text{oui, en théorie}}.
\item Les \res{\text{frottements de l'air}} freinent le jet, qui se disperse en gouttes encore plus freinées ; le vent le
dévie. En pratique, le jet n'atteint pas \qty{20}{m} : les pompiers utilisent une échelle.
\end{questions}
\end{corrige}
\end{blocLivre}
\begin{blocLivre}{cor-17-14}
\begin{corrige}{14}{La montgolfière}
\demarche{Le système décolle si la poussée, due à l'air extérieur, dépasse le poids total, air chaud compris.}
\begin{questions}
\item Poids total $P=(m_\mathrm{s}+m_\mathrm{p}+\rho_\mathrm{int}\,V)\,g$, vers le bas ; poussée $\Pi=\rho_\mathrm{ext}\,V\,g$, vers le haut.
Décollage : $\res{\Pi>P}$, soit $(\rho_\mathrm{ext}-\rho_\mathrm{int})\,V>m_\mathrm{s}+m_\mathrm{p}$.
\item $\rho_\mathrm{ext}=\dfrac{353}{288}=\res{\qty{1,23}{kg.m^{-3}}}$ et $\rho_\mathrm{int}=\dfrac{353}{373}=\res{\qty{0,946}{kg.m^{-3}}}$ ;
$\Pi=\res{\qty{3,01e4}{N}}$ et le poids de l'air chaud vaut \res{\qty{2,32e4}{N}}.
\item $(\rho_\mathrm{ext}-\rho_\mathrm{int})\,V=\qty{698}{kg}$ : la charge utile vaut $698-400=\qty{298}{kg}$, de quoi emporter
\res{\text{3 personnes}} (4 personnes pèsent \qty{320}{kg}).
\item À \qty{30}{\degreeCelsius}, $\rho_\mathrm{ext}=\qty{1,165}{kg.m^{-3}}$ : la charge utile tombe à \qty{147}{kg}, soit
\res{\text{une seule personne}}. L'air frais du matin et du soir est plus dense : la poussée est plus grande. Chauffer
davantage aide peu, l'enveloppe ne supportant guère plus de \qty{120}{\degreeCelsius}.
\end{questions}
\end{corrige}
\end{blocLivre}
\begin{blocLivre}{cor-17-15}
\begin{corrige}{15}{La centrale de haute chute}
\demarche{Bernoulli entre la surface du lac ($p_\mathrm{atm}$, $v\approx0$, $z=h$) et l'entrée de la turbine ($z=0$).}
\begin{questions}
\item $v=\dfrac{D_V}{\pi\,R^2}=\dfrac{20}{\pi\times1{,}0^2}=\res{\qty{6,37}{m.s^{-1}}}$.
\item $p=p_\mathrm{atm}+\rho\,g\,h-\frac12\,\rho\,v^2=\num{1.013e5}+\num{5.886e6}-\num{2.0e4}=\res{\qty{5,97e6}{Pa}}$, environ 60~bar. Le
terme $\frac12\,\rho\,v^2$ ne représente que \res{\qty{0,3}{\percent}} de $\rho\,g\,h$ : l'énergie arrive à la turbine surtout sous forme
de pression, d'où des conduites très épaisses en bas.
\item $P=1000\times9{,}81\times600\times20=\res{\qty{118}{MW}}$, et $P_\mathrm{é}=0{,}90\times P=\res{\qty{106}{MW}}$.
\item $E=\qty{106}{MW}\times\qty{3000}{h}=\qty{3,18e5}{MWh}$, soit $\dfrac{\num{3.18e5}}{4{,}5}\approx\res{\num{7.1e4}}$ foyers.
\end{questions}
\end{corrige}
\end{blocLivre}
